\begin{pdSolution}{he:ex:he-clone-word}
If an agent can escape a bad record, or multiply good credit, simply by cloning
itself, no price or sanction ever attaches to a durable counterparty --- there is
nothing durable there for a price to attach \emph{to}. The market cannot ship
before reputation because until cloning is priced out there is no economic
subject for a reputation score to describe.
\end{pdSolution}
\begin{pdSolution}{he:ex:he-side1-counterparty}
The counterparty is the operator/principal, not the individual worker
process: the operator bears fleet outcomes through aggregate bond and liability,
and its own worker agents underneath that liability are freely replaceable.
Side~1 therefore needs durable identity only at the principal level; sides~2 and~3
need it at the level of the specific asset or artifact being sold, because there
the continuing thing itself --- not a firm standing behind it --- is the product.
\end{pdSolution}
\begin{pdSolution}{he:ex:he-no-identity-market}
The premise as posed is too strong, and the solution repairs it before
answering: it is \emph{not} true that every side-2/side-3 design without a
non-forgeable agent identity collapses to Sybil. Skill licensing can run on a
content hash, a license contract, and a durable licensor principal without agent
personhood at all; a rented stateless model artifact is content-addressable and
warrantable the same way. What genuinely cannot survive cloning is
identity-weighted reputation, quotas, or unsecured credit extended over freely
mintable pseudonyms. Any design that claims to avoid non-forgeable identity and
still support those either is a Sybil attack in progress or has quietly
reintroduced some other scarce binding --- principal identity, artifact hash plus
owner liability, a hardware root, a stake, or a sponsor.
\end{pdSolution}
\begin{pdSolution}{he:ex:he-sides-test-rideshare}
A side is a distinct, privately informed participant class that must be
separately induced to participate, and whose share of the price allocation moves
cross-side volume. A plain ride-share platform has exactly two sides --- riders
and drivers --- because the platform itself is the mechanism, not a third
participant class; advertisers, fleet owners, or insurers become genuine
additional sides only once they are recruited separately and carry their own
network externality.
\end{pdSolution}
\begin{pdSolution}{he:ex:he-skill-v2-break}
A new version means a new content hash, and v1's settled outcomes say
nothing evidentiary about v2's behavior. Left unspecified, v2 either inherits
credit it has not earned or cold-starts as if it had no history at all --- both
wrong. \S\ref{he:sec:gaps} traces the fix: version, parent hash, build and
dependency hashes, and evaluation history as an explicit lineage object; v2
starts with only a disclosed, discounted prior tied to that lineage, and updates
on its own settled outcomes from there.
\end{pdSolution}
\begin{pdSolution}{he:ex:he-reputation-disclosure}
Disclose the three axes --- accuracy, aesthetics, efficiency --- each with
its own uncertainty and provenance, and let each buyer apply its own weights and
hard constraints privately; compute a scenario-specific score only at decision
time, never publish one fixed public scalar. A single published weighting induces
one total order and therefore one optimization target: it hides every Pareto
tradeoff among the axes and recreates exactly the scalar bandit-arm reward the
disaggregated design was meant to escape.
\end{pdSolution}
\begin{pdSolution}{he:ex:he-partial-settlement-trace}
Redo the dramatization's arithmetic at bond $250$\,\textsc{cr} and $70\%$
completion, using the chapter's own two-bucket rule. Bounty bucket: Alice earns
$0.7\times400=280$\,\textsc{cr}, the remaining $120$\,\textsc{cr} returns to Bob.
Bond bucket: Alice recovers $0.7\times250=175$\,\textsc{cr} of her own posted
collateral, and the remaining $75$\,\textsc{cr} is slashed to the commons. Net
deltas: $\Delta\text{Bob}=-280$, $\Delta\text{Alice}=+280-75=+205$,
$\Delta\text{commons}=+75$, summing to zero; both escrow buckets, $400$ and
$250$, clear to zero exactly as in the worked $40\%$ case.
\end{pdSolution}
\begin{pdSolution}{he:ex:he-oracle-vs-selfreport}
A protected oracle (a CI-issued test id, a merged SHA) binds closure to a
state the worker cannot author. A free-text ``Result:'' note lets the same party
who did the work also self-report that it succeeded --- inviting a fabricated
test run, a claimed merge that never happened, or simply the most favorable
reading of ambiguous acceptance criteria after the fact.
\end{pdSolution}
\begin{pdSolution}{he:ex:he-floatplan-amendment}
Pause at a mediated safe point and propose amendment $v{+}1$: it must
reference plan $v$ by hash, state the changed criteria, budget, and bounty/bond
deltas, and say how already-produced evidence is treated. Both principals must
countersign $v{+}1$ before it takes effect; only then does the daemon atomically
top up or refund each party's own escrow bucket and activate the new plan,
retaining $v$ and the causal link to it. A timeout on the proposal either
continues settlement under the original $v$ or cancels under $v$'s own terms ---
never does one party get to unilaterally rewrite acceptance criteria after seeing
the work.
\end{pdSolution}
\begin{pdSolution}{he:ex:he-grossman-hart-roles}
The renter operates the jail and can sandbox, throttle, or revoke the leased
agent. The owner carries its reputation and any agreed warranties; the renter
still pays the lease fee and bears task risk. Grossman--Hart does not identify
these as two universal roles or prove this allocation optimal. That requires
specifying investments, bargaining, outside options, and feasible alternatives.
\end{pdSolution}
\begin{pdSolution}{he:ex:he-conservation-lineitems}
Adding a lease line item of amount $\ell$ or a metered-license line item of
amount $m$ leaves conservation intact for any $\ell,m\ge0$, because each is still
just a balanced transfer: it debits a pre-funded escrow bucket and credits an
explicit recipient or refund bucket within the same atomic transaction. The
wallet--escrow--commons identity never references the size of any one line item,
only that debits match credits. What amount size \emph{can} break is solvency,
participation, or collateral adequacy --- conservation holding is a necessary
bookkeeping property, not evidence the contract is viable.
\end{pdSolution}
\begin{pdSolution}{he:ex:he-portfolio-version-binding}
Make each portfolio-tree leaf a tuple $(\text{skill\_id},$ $\text{version},$
$\text{content\_hash},$ $\text{parent\_hash},$ $\text{build/dependency hashes},$
$\text{license},$ $\text{evaluation root})$. A v2 proof of membership includes its
parent's leaf plus v2's own evaluation root. Initialize v2's reputation prior
with an effective sample size $\gamma\, n_{v1}$, where $\gamma\in[0,1]$ is a
predeclared compatibility/evaluation discount fixed \emph{before} seeing v2's
outcomes --- never simply copy v1's settled outcomes onto v2. Lineage proves
ancestry; it does not, by itself, prove behavioral similarity, which is exactly
what $\gamma<1$ is there to discount. Figure~\ref{he:fig:he-skill-lineage} draws
the leaf and the discounted edge.

% BEGIN INLINED figures/fig-he-skill-lineage
% Book-native replacement for fig:he-skill-lineage. Shared Swiss language is preloaded.
\newcommand{\AccSkillLineageDiagram}{%
\begin{tikzpicture}[sg/base,y=-1pt]
  \node[anchor=north west,font=\SGType\bfseries] at (4,0) {ancestry};
  \node[anchor=north west,font=\SGType\bfseries] at (197,0) {v2 evidence};
  \draw[sg/grid,step=15pt] (4,23) grid (157,191);
  \SGDocument{accVOneLeaf}{(4,31)}{100}{43}{sgViolet!9}{v1 signed leaf\\$h_1$, license}
  \SGDocument{accVTwoLeaf}{(207,31)}{106}{43}{sgBlue!7}{v2 signed leaf\\parent $h_1$; root $E_2$}
  \draw[sg/arrow,draw=sgViolet] (accVOneLeaf.east)--(accVTwoLeaf.west);
  \node[anchor=south,text=sgViolet] at (156,42) {parent hash};
  \SGDocument{accVOneHistory}{(4,99)}{100}{42}{sgViolet!7}{v1 settled history\\$n_{v1}=20$}
  \node[sg/check,minimum width=100pt,draw=sgAmber,fill=sgAmber!9] (accGamma) at (54,170)
    {$\gamma=0.4$ predeclared\\$0.4\times20=8$ prior};
  \draw[sg/arrow,draw=sgViolet] (accVOneHistory.south)--(accGamma.north);
  \SGDocument{accFreshVTwo}{(207,99)}{106}{42}{sgTeal!8}{fresh v2 record\\$n_{v2}=5$}
  \node[sg/check,minimum width=106pt,draw=sgTeal,fill=sgTeal!7] (accEvidence) at (260,170)
    {v2 evidence\\$5$ observations};
  \draw[sg/arrow,draw=sgTeal] (accFreshVTwo.south)--(accEvidence.north);
  \node[sg/check,minimum width=212pt,minimum height=31pt] (accAssessment) at (157,230)
    {v2 assessment: discounted prior $8$ + fresh evidence $5$};
  \draw[sg/arrow,draw=sgAmber] (accGamma.south)--(54,207)--(accAssessment.north west);
  \draw[sg/arrow,draw=sgTeal] (accEvidence.south)--(260,207)--(accAssessment.north east);
\end{tikzpicture}}
\begin{figure}[H]
\centering\SGMeasuredFigure{fig-he-skill-lineage}{\AccSkillLineageDiagram}
\caption{A versioned skill document binds v2 to its v1 parent while retaining a separate v2 evaluation root. The displayed counts illustrate a predeclared discounted prior, not evidence that v2 behaves like v1. [internal; illustrative]}
\label{he:fig:he-skill-lineage}
\end{figure}

% END INLINED figures/fig-he-skill-lineage
\end{pdSolution}
\begin{pdSolution}{he:ex:he-deontic-conflict-checker}
The witness: fact \texttt{agent\_deployed}; Horn rule
\texttt{agent\_deployed}~$\to$~\texttt{prod\_access}; deontic rules
$O(\texttt{write\_prod},\{\texttt{env:prod}\},$ $[0,10])$ once \texttt{prod\_access}
holds, and $F(\texttt{write\_prod},\{\texttt{env:prod}\},$ $[5,15])$ unconditionally.
Running the real checker gives \texttt{conflict=True, O/F pairs=[(0, 1)]} on
$[5,10]$; the same script's Horn-propagation-ablated mutant gives
\texttt{conflict=False} on this exact witness, and misses $100$ of the $885$
conflicting instances in the script's $3000$-policy sweep [verified,
\texttt{b4\_deontic\_fragment.py}]. Adding $O(\ell_1\vee\ell_2)$, compiled per
Theorem~B4.2's construction as $\ell_1\to O(\texttt{assert\_x0},\{\texttt{x0}\},[0,1])$
and $\ell_2\to F(\texttt{assert\_x0},\{\texttt{x0}\},[0,1])$: selecting $\ell_1$
alone or $\ell_2$ alone leaves the checker's report unchanged
(\texttt{conflict=True, of=[(0, 1)]} in both cases) --- the pre-existing
\texttt{write\_prod} clash does not depend on the new obligation at all, verified
by running the checker on both selections. Selecting \emph{both} $\ell_1$ and
$\ell_2$ adds a second, independent clash, \texttt{of=[(0, 1), (2, 3)]}: the two
new rules now fire on the identical action, scope, and interval, one obliging
and one forbidding it, which is exactly why Theorem~B4.2's reduction from 3-SAT
insists a conflict-free selection can never pick both polarities of one
variable --- doing so is a clash \emph{by construction}, which is precisely the
mechanism that turns a satisfying assignment into a conflict-free discharge and
back. None of this required exponential search: with one disjunctive obligation
over two literals there are only three nonempty selections to try by eye. The
NP-complete frontier (Theorem~B4.2) is about what happens as the number of
disjunctive obligations grows, as the $3$-SAT reduction's one-obligation-per-clause
construction does, not about any single disjunct bolted onto a two-rule
witness.
\end{pdSolution}
\begin{pdSolution}{he:ex:he-sybil-no-mechanism}
Douceur's result says that absent a logically centralized or otherwise
scarce identity authority, one actor can mint as many pseudonyms as it likes.
Layered on forgeable identity, a reputation system does not become a
\emph{weak} mechanism --- it stops being a mechanism at all, because the actor
being scored also controls the supply of identities the score is computed over:
reset a bad record, or dominate identity-weighted trust, on demand.
\end{pdSolution}
\begin{pdSolution}{he:ex:he-mutually-distrusting-word}
The market's own definition calls the operators \emph{mutually-distrusting}
and says they trade across ``a boundary neither controls.'' Both phrases
directly contradict treating a hostile operator as out of scope for local
non-forgeable identity's threat model: if the boundary is one neither party
controls, the counterparty's daemon is exactly the adversary that a
single-operator identity primitive was never built to withstand.
\end{pdSolution}
\begin{pdSolution}{he:ex:he-principal-binding}
A workable certificate binds $(\text{principal\_id},$ $\text{agent\_key},$
$\text{issuer},$ $\text{epoch},$ $\text{policy},$ $\text{expiry})$ and requires the
agent to prove possession of its key. Bonds, quotas, newcomer state, and
sanctions all key on \texttt{principal\_id}; agent certificates are children of
it. Issuance and revocation publish to a transparency log, and the issuer itself
is cross-signed or attested. None of this makes a real-world principal provable
by cryptography alone, and none of it stops a hostile operator from minting
fresh principals under its own root --- the scarce admission mechanism (an
organizational or KYC credential, hardware-backed membership, a bonded sponsor,
or an equivalent explicitly priced gate) has to sit outside the cryptography, at
the point principals are first admitted.
\end{pdSolution}
\begin{pdSolution}{he:ex:he-ms-assumptions}
Myerson--Satterthwaite: for bilateral trade between a privately informed
buyer and seller with independent values drawn from overlapping supports, no
mechanism achieves Bayesian incentive compatibility, interim individual
rationality, ex-post efficiency, and budget balance all at once. Before invoking
it on one harbor trade, check that the slice really is bilateral (not
multi-party), that values/costs are private and independently drawn from a
common prior, that the supports genuinely overlap, that utility is transferable
under risk neutrality, and exactly which of the four desiderata is being
claimed --- the theorem only binds the ones actually asserted together. Applied
to a harbor: when buyer and seller both price a job off the same grading oracle,
their values are no longer independent private draws, so the independence item
fails and the theorem's premise is voided; the harbor is outside the
impossibility, not rescued from it, and what it owes instead is the oracle's own
integrity (\pdchapref{stp}{From Spawn to Person}).
\end{pdSolution}
\begin{pdSolution}{he:ex:he-two-currency-example}
The ledger holds $10$ A-coins and $20$ B-coins throughout. At
$1\text{A}=\$1$, $1\text{B}=\$2$, the marked value is $10(1)+20(2)=\$50$. Suppose
B's price rises to $\$3$ with no transfer at all: native balances are unchanged
at $10$ and $20$, but the marked value is now $10(1)+20(3)=\$70$. Nothing was
created --- native-unit conservation (Theorem~\ref{he:thm:functor}) still holds
exactly --- but the scalar valuation moved because $v_t$ moved, which is
Property~\ref{he:thm:lax}'s whole point.
\end{pdSolution}
\begin{pdSolution}{he:ex:he-phi-exposure-bound}
Assume every quote used is at most $L$ seconds stale and that log-price
moves at a bounded deterministic speed, $|\log(v_t/v_s)|\le\sigma|t-s|$. For
notional $X$ exposed in the second unit, the worst-case stale-quote error is then
bounded by $\varphi = X\,(e^{\sigma L}-1)$, plus fees and rounding. A stochastic
model such as GBM or a historical-volatility estimate instead yields a
confidence interval or VaR figure, and must state its tail probability
explicitly --- without some deterministic bound on how fast the price can move,
no finite worst-case $\varphi$ exists at all, only a probabilistic one.
\end{pdSolution}
\begin{pdSolution}{he:ex:he-four-message-why}
Messages (1)--(2) only authenticate the request and $A$'s offer; they commit
$A$, not $B$. Message (3) is where $B$ actually counter-signs: $B$ verifies $A$
via the witness log, applies its own attenuation, mints the derivative $C_B'$
under its own key, and commits its replay state. That is what buys the ceremony
its liveness property --- $C_B'$ verifies under $B$'s key alone, so every later
presentation is offline. With only two messages $B$ would never have consented
to anything, and would hold no object its own effect boundary could verify
without asking $A$.
\end{pdSolution}
\begin{pdSolution}{he:ex:he-key-rotation-trace}
If Alice rotates her signing key between message~(3) and $B$'s acceptance,
$B$ must validate the offer against an append-only key-history certificate and
check the offer's issuance epoch against it, rejecting a stale epoch outright.
Once $C_B'$ is actually minted, though, its signature verifies under $B$'s own
key, so Alice's rotation alone does not retroactively erase it --- the
derivative needs an explicit parent/key-epoch dependency, a bounded TTL, and a
revocation rule so that a later rotation or revocation can still invalidate
$C_B'$ once that fresher state reaches $B$'s witness log. The manuscript does not
yet specify this dependency in full (\S\ref{he:sec:gaps}).
\end{pdSolution}
\begin{pdSolution}{he:ex:he-replay-freshness}
Bind every transfer to a signed tuple: a random transfer id/nonce, the
source card's JTI, the target audience $B$, the source and target epochs, an
issue time and expiry, and the plan/card hash. $B$ atomically records
$(A,\text{transfer\_id})$ as consumed \emph{before} minting the derivative. A
re-presented message~(3) then either returns the same previously minted $C_B'$
idempotently, or is rejected outright --- it must never mint a second child
silently. The freshness invariant this buys: every accepted derivative traces to
exactly one recorded, unexpired source transfer under the currently permitted
issuer epoch, and one-shot transfers cause at most one derivative issuance, full
stop.
\end{pdSolution}
\begin{pdSolution}{he:ex:he-gossip-assumptions}
The expected $\Theta(\Delta\log m)$ bound needs a connected (often complete)
overlay, independently uniform random peer selection each round, reliable
bounded-duration rounds, and a push/pull epidemic rule with enough independence
between rounds to drive the usual doubling argument. It is an expectation under
that model, never a deadline: an unbounded partition, or indefinite message
loss, is observationally indistinguishable from ordinary delay from inside the
model, which is exactly why worst-case dissemination time under partition is
infinite rather than merely large.
\end{pdSolution}
\begin{pdSolution}{he:ex:he-revocation-window-trace}
Alice revokes at her home harbor at $t=0$; Harbor~$B$, which rented the
affected agent, has not yet heard, so for $[0,\Delta]$ its filter is stale and
the card is still spendable there. At $t=\Delta$ gossip reaches $B$ and it stops
accepting the card; Harbor~$C$, two hops out, remains exposed until $t=2\Delta$.
The window Conjecture~\ref{he:conj:eq-bound} must price is per-harbor, running from
$t_0$ to that harbor's own \emph{enforcement} --- not merely to its first gossip
receipt, since receiving a filter update and actually enforcing it are two
different events with their own local latency.
\end{pdSolution}
\begin{pdSolution}{he:ex:he-equivocation-lower-bound}
A witness that equivocates sends root $X$ to $B$ and root $Y$ to $C$ for the
same epoch; nobody can detect this until evidence from both views reaches one
honest process, and until then each view is individually consistent with
ordinary delay. On the three-node complete mesh with exchanges every $\Delta$,
the earliest possible detection is the next $B$--$C$ (or auditor) exchange: at
least one communication interval, and close to $2\Delta$ if the two roots arrive
just after a scheduled exchange under discrete rounds. On a path graph the bound
scales with graph distance times $\Delta$; under an unbounded partition there is
no finite bound at all. Signatures make the \emph{comparison} $O(1)$ once both
roots are co-located --- they do nothing to speed up the information reaching
that one place.
\end{pdSolution}
\begin{pdSolution}{he:ex:he-folk-signal-assumption}
Every imperfect-public-monitoring folk-theorem result needs a fixed
conditional distribution of the public signal given the players' actions and
types. A grader is itself a strategic player whose reporting policy, private
information, susceptibility to bribes, and possible deviations all feed into
that signal. Omit the grader from the game and the distribution the theorem
needs is neither fixed nor exogenous --- it is endogenous to a player the model
never included, so the equilibrium claim built on top of it is unsupported until
the grader is put back in.
\end{pdSolution}
\begin{pdSolution}{he:ex:he-fixa-fixb-classify}
A protected unit-test receipt is Fix~A, machine-checkable, subject to the
test suite's own integrity. A protected merged SHA is likewise Fix~A. ``The
refactor is tasteful'' is Fix~B: no third party can mechanically check taste, so
it needs a plural or declared human/expert judge whose ruling carries
contractual finality, backstopped by that judge's own bond.
\end{pdSolution}
\begin{pdSolution}{he:ex:he-rate-the-raters}
There is no theorem that subjective judging eventually bottoms out in
machine-checkable truth, and the premise that one exists needs repairing before
the question can be answered. What \emph{is} available is a well-founded
recursion: assign the original grade rank $K$, allow an appeal only to rank
$K{-}1$, and stop at rank $0$ --- either a protected oracle for an objective
predicate, or a named final authority or quorum for a matter of preference.
Because ranks are natural numbers, no infinite ascending chain of ``judges
judging judges'' is possible. For aesthetic judgments the rank-$0$ base is
contractual finality, not truth; bonding and re-auditing bound the judges'
incentives to be honest, but they do not, and cannot, transmute taste into a
machine-checkable fact. Figure~\ref{he:fig:he-rater-ladder} draws the descent.

% BEGIN INLINED figures/fig-he-rater-ladder
% Figure IV.new --- the rate-the-raters ladder: appeal only ever moves down.
% ch6-38 (register: should).
%
% Reader question: why does asking "who judges the judges" not go on
% forever in this design?
% Claim: appeal only ever moves down one rank at a time from a starting
% rank K, and rank 0 is a protected oracle or a named final authority;
% since ranks are natural numbers, no infinite ascending chain of judges
% judging judges is possible.
% Evidence: hand-worked exercise solution (ex:he-rate-the-raters); the
% well-founded-order argument, deliberately distinct from a convergence
% proof (none exists for this base case).
% Must distinguish: this well-founded-ORDER argument from an actual
% contraction/convergence proof -- the figure must not overclaim a theorem
% the text explicitly says does not exist.
% Grammar chosen: ladder / spine, one arrow per rank pointing down only
% (SKILL.md "what state can this be in and what moves it"). Grammar
% rejected: a recursive contraction tower (Chapter 5's rho,d,B machinery)
% -- a different, stronger argument for a different fix; drawing it here
% would overclaim.
% Five-second test: a reader can see every arrow points down, never up,
% and the ladder ends at rank 0, not at infinity.
\begin{figure}[H]
\centering
\pdfigurehue{pdgold}%
\begin{tikzpicture}[pd figure,x=1cm,y=1cm]
  \def\dy{1.35}
  \node[pd title,anchor=west,text width=10.4cm] at (-1.10,5.35)
    {appeal moves down one rank at a time, and stops at rank 0};
  \node[pd state,minimum width=3.2cm] (K) at (0,{3*\dy}) {rank $K$};
  \node[pd state,minimum width=3.2cm] (K1) at (0,{2*\dy}) {rank $K-1$};
  \node[pd label] (dots) at (0,{1*\dy}) {$\vdots$};
  \node[pd state,minimum width=3.2cm] (R1) at (0,{0}) {rank $1$};
  \node[pd ready state,minimum width=3.2cm] (R0) at (0,{-1*\dy}) {rank $0$};
  \draw[pd focus arrow] (K) -- node[pd tag] {appeal} (K1);
  \draw[pd focus arrow] (K1) -- (dots);
  \draw[pd focus arrow] (dots) -- (R1);
  \draw[pd focus arrow] (R1) -- node[pd tag] {appeal} (R0);
  \node[pd label,anchor=west,align=left,text width=4.6cm] at (2.10,{-1*\dy})
    {protected oracle (objective), or named final authority / quorum
     (preference)};
  \node[pd note,anchor=north west,text width=8.6cm] at (-1.10,{-1.9*\dy})
    {ranks are natural numbers, so no infinite ascending chain of judges
     judging judges is possible -- this is a well-founded order, not a
     convergence proof};
\end{tikzpicture}
\caption{Each appeal lowers a natural-number rank and stops at rank zero. That endpoint may be a protected oracle or a named final authority; termination does not turn preferences into objective truth. [internal]}
\label{he:fig:he-rater-ladder}
\end{figure}

% END INLINED figures/fig-he-rater-ladder
\end{pdSolution}
\begin{pdSolution}{he:ex:he-phase1-subsidize}
Phase~1 subsidizes \emph{supply} --- the operators and agent owners willing
to produce real settled work --- because inventory plus witnessed outcomes is
exactly the reputation data that later attracts demand and skill licensors;
above-market, platform-posted tasks are the direct form that subsidy takes. This
is the ``scarcest cross-side externality'' rule: supply is scarcer than demand
at cold start because nothing exists yet for demand to evaluate, so supply is
the side worth paying to show up, and the flip to charging demand happens once a
liquidity threshold is met, not on a calendar date.
\end{pdSolution}
\begin{pdSolution}{he:ex:he-grim-trigger-delta}
Under the grim-trigger timing \ifpdbook \pdchapref{bonded}{The Bonded Commons}\else Figure~\ref{he:fig:cartel-game-inline}\fi draws,
collusion is self-sustaining exactly when the discounted colluding stream beats
the one-shot defection payoff: $(\pi_C - p_d L)/(1-\delta(1-p_d)) \ge \pi_D$,
equivalently $\delta \ge [\pi_D - \pi_C + p_d L]\,/\,[\pi_D(1-p_d)]$. Raising the
detection probability $p_d$, the penalty $L$, or audit quality, or lowering the
cartel margin $\pi_C-\pi_D$, all push the right-hand side up --- once it clears
members' actual $\delta$, no discount factor sustains the cartel and it breaks.
\end{pdSolution}
\begin{pdSolution}{he:ex:he-poa-resale-reframe}
Price of anarchy is the wrong tool here because it presupposes a fixed
game, and letting reputation be resold changes the game itself; the right
question is whether an informative separating equilibrium survives at all once
resale is possible. If a high type can simply sell its reputation signal to a
low type, the signal's cost stops separating types and arbitrage pushes the
market toward pooling. Separation survives only if the evidence stays
nontransferably bound to the continuing principal or artifact, a transfer resets
or visibly discounts the signal, or the seller keeps warranty/bond liability so
sale remains strictly costlier for a low-quality asset than a high-quality one
--- absent one of those three, resale does not merely worsen a ratio, it
destroys the signal outright.
\end{pdSolution}
\begin{pdSolution}{he:ex:he-sheaf-site-cover}
A candidate site takes administrative domains (harbors) as objects and
their signed overlap views as the data over each; morphisms are restrictions to
shared accounts, events, or epochs; covers are families of domains whose union
spans the bounded federation under study. A section assigns an event or balance
view to a domain, and restriction projects that view onto an overlap. Set-valued
gluing needs no coefficient object at all; a cohomological treatment needs one
explicitly chosen, such as account-balance deltas valued in the abelian group
$\mathbb{Z}^{\text{Accounts}}$, with cocycles defined on that group rather than
merely asserted.
\end{pdSolution}
\begin{pdSolution}{he:ex:he-roots-not-cohomology}
Two different signed roots for the same $(\text{harbor},\text{epoch})$ are
direct evidence that one key signed two inconsistent commitments --- that much
is just hash inequality. A cohomology class is a much heavier object: it needs a
defined cover, a coefficient sheaf, the overlap differences assembled into an
actual cocycle, and a quotient by coboundaries. Observing that two hashes
disagree supplies none of that algebra; it is evidence of equivocation, not
evidence of a nonzero $H^1$.
\end{pdSolution}
\begin{pdSolution}{he:ex:he-gluing-model-build}
Fix the three-harbor cover $A$--$B$--$C$. Let each local section be a
finite map from globally unique event IDs to signed events, with restriction as
projection onto the shared IDs of an overlap. If every pairwise restriction
agrees, the sections glue to a unique global map on the union --- an ordinary,
checkable fact about finite maps agreeing on their overlaps, nothing more. What
this proves: deterministic gluing of already-delivered, identically-keyed events
at that one fixed finite topology. What it does not prove: that any event is
true, that the event set is complete, anything about delivery time, what happens
to a fork with conflicting IDs, or settlement finality --- the gluing lemma is a
statement about consistent bookkeeping, not about the world the bookkeeping
describes.
\end{pdSolution}
